% Partie : théorie de la démonstration
% Chapitre : normalisation
% Section : correspondance des preuves normales et des séquents

\documentclass[../../../include/open-logic-section]{subfiles}

\begin{document}

\olfileid[fr]{pt}{nor}{tr}

\olsection{Correspondance entre les \usetoken{p}{proof} normales de \Log{N2i} et \Log{G2i}}

Les preuves de \Log{N2i} et celles de \Log{G2i}
peuvent être traduites les unes dans les autres.
En particulier, une preuve \emph{sans coupure}
de \Log{G2i} donne une preuve normale de
\Log{N2i}.

Rappelons la correspondance des contextes.
Un ensemble~$\Gamma'$ de formules étiquetées
\emph{correspond} à un multiensemble~$\Gamma$
de formules sans étiquette si, pour chaque
formule~$!A$, le nombre d'éléments de
$\Gamma'$ de la forme $x : !A$ est inférieur
ou égal à la multiplicité de~$!A$ dans~$\Gamma$.

\begin{cor}
Si $\Gamma \Sequent \Delta$ admet une preuve
sans coupure dans \Log{G2i}, sa traduction
est une preuve normale~$\delta$ de
$\Gamma' \Sequent !C$ dans \Log{N2i}.
Le contexte~$\Gamma'$ correspond au
multiensemble~$\Gamma$. Si le conséquent
est un singleton, sa formule est~$!C$.
Si le conséquent est vide, $!C$ est l'absurde.
Chaque formule~$!A$ a dans~$\Gamma'$ au plus
autant de classes étiquetées que
d'occurrences dans~$\Gamma$.
\end{cor}

\begin{proof}
Examinons la traduction de la
\olref[nat][tgi]{prop:G2i-to-N2i}.
Les introductions sont ajoutées à la fin
des sous-preuves. Les prémisses majeures
des éliminations proviennent des axiomes
des classes désignées, y compris lorsque
l'élimination réunit des branches déjà
traduites. Renommer ou fusionner des
classes ne change pas les inférences.
Ainsi, sans traduction de \Cut, aucune
élimination ne prend pour majeure la
conclusion d'une introduction, de~\FalseInt,
ni la dernière occurrence d'un segment
de longueur supérieure à un.
\end{proof}

La traduction de \Cut{} dans la
\olref[nat][tgi]{prop:G2i-to-N2i} peut,
en revanche, créer une coupure dans
la preuve~$\delta$ de \Log{N2i}.
Si une preuve de \Log{G2i} se termine par
une coupure dont les sous-preuves
$\pi_1$ et~$\pi_2$ concluent
$\Gamma_1 \Sequent !A$ et
$!A, \Gamma_2 \Sequent \Delta_2$,
on greffe la traduction~$\delta_1$
de~$\pi_1$ aux axiomes
$x:!A \Sequent !A$ de la classe
désignée dans la traduction~$\delta_2$
de~$\pi_2$. Si~$\delta_1$ se termine
par une introduction de formule
principale~$!A$, et si l'axiome
$x:!A \Sequent !A$ est une prémisse
majeure d'élimination, le greffage
crée un détour.
Les étiquettes et les constantes propres
sont préalablement rendues fraîches.

Dans l'autre sens, la traduction générale
de \Log{N2i} vers \Log{G2i} donnée dans la
\olref[nat][tni]{prop:N2-to-G2} emploie
des coupures pour \Elim{\land},
\Elim{\lif}, \Elim{\lnot} et
\Elim{\lforall}. Si la preuve initiale
est normale, on peut les éviter.

Une \emph{branche} d'une preuve~$\delta$
de \Log{N2i} est une suite $S_1$, \dots,
$S_n$ d'occurrences de séquents,
commençant par un axiome, telle que
chaque $S_i$, avant le dernier, est
une prémisse majeure d'une inférence
dont $S_{i+1}$ est la conclusion.
La branche s'arrête lorsqu'elle atteint
une prémisse mineure d'élimination.
Elle est \emph{principale} si $S_n$
est le séquent final.
Toute preuve~$\delta$ possède au moins
une branche principale : chaque règle
a une prémisse majeure, et seule
\Intro{\land} en a deux.

\begin{prop}\ollabel{prop:normal-N2i-to-G2i}
Si $\delta$ est une preuve normale
de $\Log{N2c}$ ou de $\Log{N2i}$
de~$\Gamma \Sequent !A$, il existe,
respectivement, une preuve sans coupure
de $\Log{G2c}$ ou de $\Log{G2i}$
de~$\Gamma' \Sequent \Delta$.
Le multiensemble~$\Gamma'$ est obtenu
en effaçant les étiquettes de~$\Gamma$,
avec une copie par classe d'hypothèses.
On prend $\Delta = \{!A\}$ si~$!A$
n'est pas~$\lfalse$, et
$\Delta = \emptyset$ sinon.
\end{prop}

\begin{proof}
Nous donnons d'abord la transformation
directe pour la version intuitionniste.
La récurrence porte sur le nombre
d'inférences de la preuve~$\delta$
de~$\Gamma \Sequent !A$.
Sans inférence, celle-ci est un axiome
$x:!A \Sequent !A$. Sa traduction~$\pi$
est $!A \Sequent !A$, ou
$\lfalse \Sequent \ $ si
$!A \ident \lfalse$.

Dans les constructions qui suivent,
le conséquent final reste vide lorsque
la conclusion naturelle est l'absurde.
On rétablit une formule $\lfalse$ à
droite par \RightR{\Weakening} seulement
si une prémisse logique l'exige.
On n'a jamais besoin de \Cut{} pour
retirer un absurde du conséquent final.
Les contextes naturels sont des ensembles
de classes ; les contextes des séquents
sont des multiensembles. Lorsque deux
prémisses utilisent la même classe
ouverte, leurs traductions donnent deux
copies de sa formule : on les contracte
pour en garder une par classe.
Des classes distinctes de même formule
demeurent, elles, des copies distinctes.

Supposons que la dernière inférence
soit une règle~$R$.
\begin{enumerate}
  \item Si $R$ est une introduction
  ou~\FalseInt, la traduction suit celle de la
  \olref[nat][tni]{prop:N2-to-G2},
  sans coupure. Une décharge vacante
  demande, au besoin, un affaiblissement
  gauche. Pour une introduction dont
  une prémisse est l'absurde, on rétablit
  cette formule à droite si le schéma
  l'exige. L'introduction de négation
  utilise directement le conséquent vide.
  Pour~\FalseInt, on conserve ce
  conséquent vide si la conclusion est
  encore l'absurde ; sinon, on applique
  l'affaiblissement droit avec la formule
  finale. Les contextes communs sont
  contractés comme indiqué.

  Si $R$ est une élimination, nous
  utiliserons les observations suivantes.
  \begin{enumerate}
    \item Aucune branche principale
    ne contient d'introduction ni
    de~\FalseInt : une telle règle,
    suivie sur cette branche d'une
    élimination, formerait un segment
    de coupure.
    \item La branche principale est
    donc unique. Une bifurcation de
    branches principales ne pourrait
    provenir que de \Intro{\land}.
    \item Si elle passe par la prémisse
    majeure de \Elim{\lor} ou de
    \Elim{\lexists}, cette règle est
    nécessairement la dernière, donc~$R$.
    Sinon, sa conclusion finirait par
    servir de prémisse majeure d'une
    autre élimination et formerait un
    segment de coupure.
    \item Les éliminations autres que
    \Elim{\lor} et \Elim{\lexists}
    ne déchargent aucune hypothèse.
    Ces deux règles ne déchargent que
    des classes de leurs prémisses
    mineures. Si~$S_i$ sur la branche
    principale a pour contexte~$\Gamma_1$,
    le contexte de~$S_{i+1}$
    contient donc celui-ci :
    $\Gamma_1' \supseteq \Gamma_1$.
    \item Si l'axiome initial~$S_1$
    est $x: !C \Sequent !C$,
    son hypothèse reste ouverte ;
    en particulier, $x:!C \in \Gamma$.
  \end{enumerate}
  Nous distinguons maintenant la forme
  de~$!C$, prémisse majeure de la première
  élimination.

  Avant de remplacer cet axiome, nous
  isolons son occurrence par une étiquette
  fraîche, notée~\emph{x} dans les dessins.
  Les autres occurrences de sa classe
  initiale, si elles existent dans des
  branches mineures, gardent une autre
  étiquette. Cette séparation conserve
  la validité et le nombre d'inférences,
  puisque l'occurrence principale ne
  sera déchargée par aucune règle.
  Après la traduction, une contraction
  réunit les copies issues de cette même
  classe initiale. Lors d'un remplacement,
  tous les contextes de la continuation
  sont recalculés à partir des feuilles :
  une dépendance encore utilisée dans
  une autre branche n'est pas supprimée.

  \item Si $!C \ident !D \land !E$,
  la preuve a, pour la première projection,
  la forme
\[
  \Axiom$x:!D \land !E \fCenter !D \land !E$
  \RightLabel{\Elim{\land}[1]}
  \UnaryInf$x:!D \land !E \fCenter !D$
  \Deduce$x:!D \land !E, \Gamma_1 \fCenter !A$
  \DisplayProof
  \]
  Sur la branche principale, l'hypothèse
  $x:!D \land !E$ reste ouverte.
  Le séquent $x:!D \land !E \Sequent !D$
  se trouve sur cette branche. La nouvelle
  classe~$x:!D$ remplacera seulement
  l'occurrence principale isolée.
  Nous remplaçons la sous-preuve se
  terminant par cette projection par
  l'axiome $x:!D \Sequent !D$,
  et recalculons les contextes dans
  la continuation. Le contexte résiduel
  contient toute dépendance des autres
  branches, y compris une éventuelle
  autre occurrence de la classe initiale.
  Schématiquement, nous transformons
\[
  \bottomAlignProof
  \Axiom$x:!D \land !E \fCenter !D \land !E$
  \RightLabel{\Elim{\land}[1]}
  \UnaryInf$x:!D \land !E \fCenter !D$
  \Deduce$x:!D \land !E, \Gamma_1 \fCenter !A$
  \DisplayProof
  \quad\text{ en }\quad
  \bottomAlignProof
  \Axiom$x:!D \fCenter !D$
  \Deduce$x:!D, \Gamma_1 \fCenter !A$
  \DisplayProof
  \]
  La nouvelle preuve~$\delta_1$ est
  normale et possède une inférence
  de moins que~$\delta$.
  L'hypothèse de récurrence donne
  une preuve~$\pi_1$ sans coupure,
  de $!D, \Gamma_1' \Sequent !A$
  lorsque la conclusion n'est pas l'absurde.
  On applique \LeftR{\land} :
\[
  \AxiomC{}
  \RightLabel{$\pi_1$}
  \Deduce$!D, \Gamma_1' \fCenter !A$
  \RightLabel{\LeftR{\land}}
  \UnaryInf$!D \land !E, \Gamma_1' \fCenter !A$
  \DisplayProof
  \]
  Si $!A \ident \lfalse$, on garde
  le conséquent vide :
\[
  \AxiomC{}
  \RightLabel{$\pi_1$}
  \Deduce$!D, \Gamma_1' \fCenter $
  \RightLabel{\LeftR{\land}}
  \UnaryInf$!D \land !E, \Gamma_1' \fCenter $
  \DisplayProof
  \]
  Les éventuelles copies de la même
  classe initiale sont ensuite contractées.
  La seconde projection se traite avec
  l'autre règle \LeftR{\land}.

  \item Si $!C \ident !D \lor !E$,
  l'axiome initial est la prémisse majeure
  de $\Elim{\lor}$, qui est la dernière
  règle~$R$. La preuve a la forme
\[
  \Axiom$x:!D \lor !E \fCenter !D \lor !E$
  \AxiomC{}
  \RightLabel{$\delta_1$}
  \Deduce$y:!D, \Gamma_1 \fCenter!A$
  \AxiomC{}
  \RightLabel{$\delta_2$}
  \Deduce$z:!E, \Gamma_2 \fCenter !A$
  \RightLabel{\Elim{\lor}}
  \TrinaryInf$x:!D \lor !E, \Gamma_1, \Gamma_2 \fCenter !A$
  \DisplayProof
  \]
  La récurrence s'applique aux deux
  branches~$\delta_1$ et~$\delta_2$.
  Les contextes résiduels $\Gamma_1$ et~$\Gamma_2$
  s'obtiennent après la décharge éventuelle
  de $y:!D$ et de $z:!E$, respectivement.
  Les traductions~$\pi_1$
  et~$\pi_2$ sont combinées par la règle
  intuitionniste $\LeftR{\lor}$ à
  \emph{conséquent partagé} :
\[
  \AxiomC{}
  \RightLabel{$\pi_1$}
  \Deduce$!D, \Gamma_1' \fCenter !A$
  \AxiomC{}
  \RightLabel{$\pi_2$}
  \Deduce$!E, \Gamma_2' \fCenter !A$
  \RightLabel{\LeftR{\lor}}
  \BinaryInf$!D \lor !E, \Gamma_1', \Gamma_2' \fCenter !A$
  \DisplayProof
  \]
  Si une classe désignée n'est pas utilisée,
  un affaiblissement gauche ajoute le
  disjoint correspondant. Lorsque~$!A$
  est l'absurde, les deux conséquents
  partagés sont vides. Par exemple,
  si les deux décharges sont vacantes :
\[
  \AxiomC{}
  \RightLabel{$\pi_1$}
  \Deduce$\Gamma_1' \fCenter $
  \RightLabel{\LeftR{\Weakening}}
  \UnaryInf$!D, \Gamma_1' \fCenter $
  \AxiomC{}
  \RightLabel{$\pi_2$}
  \Deduce$\Gamma_2' \fCenter $
  \RightLabel{\LeftR{\Weakening}}
  \UnaryInf$!E, \Gamma_2' \fCenter $
  \RightLabel{\LeftR{\lor}}
  \BinaryInf$!D \lor !E, \Gamma_1', \Gamma_2' \fCenter $
  \DisplayProof
  \]
  Aucun affaiblissement droit final
  n'est ajouté. On contracte ensuite
  seulement les copies issues d'une
  même classe ouverte.

  \item Si $!C \ident !D \lif !E$,
  la première élimination donne
\[
  \Axiom$x:!D \lif !E \fCenter !D \lif !E$
  \AxiomC{}
  \RightLabel{$\delta_1$}
  \Deduce$\Gamma_1 \fCenter!D$
  \RightLabel{\Elim{\lif}}
  \BinaryInf$x:!D \lif !E, \Gamma_1 \fCenter !E$
  \RightLabel{$\delta_2$}
  \Deduce$x:!D \lif !E, \Gamma_1, \Gamma_2 \fCenter !A$
  \DisplayProof
  \]
  Le séquent
  $x:!D \lif !E, \Gamma_1 \Sequent !E$
  se trouve sur la branche principale ;
  son contexte est $x:!D \lif !E, \Gamma_1$.
  Nous conservons la preuve~$\delta_1$
  de l'argument et remplaçons la
  sous-preuve concluant l'élimination
  par un axiome frais $x:!E \Sequent !E$.
  La nouvelle classe~$x:!E$ ne porte
  pas l'étiquette d'une hypothèse
  conservée dans une autre branche.
  Dans la continuation, notée
  $\delta_2'$, les contextes sont
  recalculés. Le contexte résiduel
  conserve toutes les
  hypothèses encore employées par ses
  autres branches, même si elles
  figuraient aussi dans~$\Gamma_1$.
  Le changement est schématisé par
\[
  \bottomAlignProof
  \Axiom$x:!D \lif !E, \Gamma_1 \fCenter !E$
  \RightLabel{$\delta_2$}
  \Deduce$x:!D \lif !E, \Gamma_1, \Gamma_2 \fCenter !A$
  \DisplayProof
  \quad\text{ en }\quad
  \bottomAlignProof
  \Axiom$x:!E \fCenter !E$
  \RightLabel{$\delta_2'$}
  \Deduce$x:!E, \Gamma_2 \fCenter !A$
  \DisplayProof
  \]
  Les deux nouvelles preuves sont
  normales et plus petites. Le nombre
  d'inférences de~$\delta$ est la somme
  de ceux de~$\delta_1$ et de
  $\delta_2'$, plus~$1$ pour l'élimination
  supprimée. Si cette élimination était
  déjà la dernière règle~$R$,
  c'est~$\delta_2'$ qui contient
  $0$ inférence ; la preuve de
  l'argument peut être non triviale.

  Les traductions~$\pi_1$ et~$\pi_2$
  concluent $\Gamma_1' \Sequent !D$
  et $!E, \Gamma_2' \Sequent !A$
  lorsque ces formules ne sont pas
  l'absurde. Elles sont combinées par
  \LeftR{\lif} :
\[
  \AxiomC{}
  \RightLabel{$\pi_1$}
  \Deduce$\Gamma_1' \fCenter !D$
  \AxiomC{}
  \RightLabel{$\pi_2$}
  \Deduce$!E, \Gamma_2' \fCenter !A$
  \RightLabel{\LeftR{\lif}}
  \BinaryInf$!D \lif !E, \Gamma_1', \Gamma_2' \fCenter !A$
  \DisplayProof
  \]
  Si $!D$ est l'absurde, on rétablit
  cette formule à droite de la première
  prémisse. Si $!A \ident \lfalse$,
  la seconde prémisse et la conclusion
  gardent un conséquent vide.
  Ainsi, lorsque les deux sont l'absurde :
\[
  \AxiomC{}
  \RightLabel{$\pi_1$}
  \Deduce$\Gamma_1' \fCenter $
  \RightLabel{\RightR{\Weakening}}
  \UnaryInf$\Gamma_1' \fCenter !D$
  \AxiomC{}
  \RightLabel{$\pi_2$}
  \Deduce$!E, \Gamma_2' \fCenter $
  \RightLabel{\LeftR{\lif}}
  \BinaryInf$!D \lif !E, \Gamma_1', \Gamma_2' \fCenter $
  \DisplayProof
  \]
  Des contractions rétablissent une
  copie par classe du contexte initial.
\end{enumerate}

Les autres éliminations se traitent de
la même manière. Pour la négation,
la première règle est \Elim{\lnot} :
on traduit sa prémisse mineure, en
rétablissant l'absurde à droite si
l'argument est lui-même l'absurde,
puis on applique \LeftR{\lnot}.
La conclusion a un conséquent vide.
Pour l'universelle, on remplace la
première élimination par l'axiome
sur son instance, traduit la continuation,
puis applique \LeftR{\lforall} avec
le même terme fermé.
Pour l'existentielle, la première règle
est aussi la dernière : on traduit sa
branche mineure, ajoute l'hypothèse
instanciée par affaiblissement gauche
en cas de décharge vacante, puis applique
\LeftR{\lexists}. La régularité et les
conditions d'élimination garantissent
la fraîcheur de sa constante propre.
Les éventuels contextes communs sont
contractés dans tous ces cas.
Ces constructions n'emploient aucune
coupure.

La partie classique de l'énoncé ne
découle pas de cette analyse de branche,
car~\FalseCl{} peut décharger une
hypothèse sur la branche principale.
Elle s'obtient séparément en combinant
la traduction générale de \Log{N2c}
vers \Log{G2c} avec coupures et le
théorème classique d'élimination des
coupures. Cette voie donne le séquent
requis sans exiger la normalité de la
preuve initiale ; elle n'affirme pas
que chaque formule de cette preuve
survit dans la preuve finale.
\end{proof}

\begin{cor}
Dans une preuve normale de \Log{N1i} ou
de \Log{N2i}, toute formule autre que
$\lfalse$ est une sous-formule de la
conclusion ou d'une hypothèse ouverte,
au sens qui admet les instances de
termes pour les sous-formules quantifiées.
\end{cor}

\begin{proof}
Par la \olref{prop:normal-N2i-to-G2i}, une
preuve normale de \Log{N2i} se traduit
en une preuve sans coupure de \Log{G2i}.
La transformation directe conserve les
formules de la preuve naturelle, à
l'exception éventuelle de l'absurde,
que le conséquent vide peut représenter.
La propriété de la sous-formule du calcul
sans coupure s'applique aux autres
formules, avec les instances requises
par les règles quantifiées.
Le passage de \Log{N1i} à \Log{N2i}
conserve les occurrences de formules
et explicite leurs hypothèses ouvertes.
\end{proof}

L'exception est nécessaire lorsque la
négation est primitive. À partir de
$!A$ et de $\lnot !A$, on obtient
$\lfalse$ par élimination de la négation,
puis un atome $!B$ distinct de~$!A$ par
\FalseInt. Pour $!A$ atomique, cette
preuve est normale et $\lfalse$ n'est
une sous-formule ni de ses hypothèses
%
%
%
%
%
%
%
%
%
%
%
%
%
%
%
%
%
%
%
%
%
ni de sa conclusion.

\end{document}
